% !TeX root = ../main.tex

% ==================================================
% A.1. 함수 객체, 이름 바인딩, callable, 호출과 반환
% ==================================================

\subsection{함수 객체, 이름 바인딩, callable, 호출과 반환}

\subsubsection{클래스 객체 모델과의 연결}

클래스 객체 모델에서 이런 코드를 살펴볼 수 있다.

\begin{dhkcode}
class Person:
    def introduce(self):
        return "Hello"


p = Person()

method = p.introduce

print(method.__self__ is p)
print(method.__func__ is Person.introduce)
\end{dhkcode}

실행 결과:

\begin{dhkzsh}
True
True
\end{dhkzsh}

우리는 이 결과를 디스크립터 프로토콜(Descriptor Protocol)을 통해 설명한다.

\zsh{Person.__dict__["introduce"]}에는 함수 객체가 저장되고,

\zsh{p.introduce}에 접근하면 함수가 인스턴스에 바인딩된 메서드 객체를 얻는다.

그런데 여기에는 아직 독립적으로 설명하지 않은 전제가 있다.

클래스 네임스페이스에 저장된 함수 객체는 애초에 어떻게 생성되었으며, Python은 그 객체를 어떻게 호출하는가?

또한 다음 표현식을 이미 사용하고 있었다.

\begin{dhkcode}
Person.introduce
Person.introduce(p)
\end{dhkcode}

두 표현식은 괄호의 유무만 다른 것처럼 보이지만, 실제로는 서로 다른 객체 연산을 나타낸다.

첫 번째는 속성 탐색으로 얻은 함수 객체를 평가하고, 두 번째는 그 함수 객체를 호출한다.

이번 학습에서는 이 차이를 출발점으로 삼을 것이다.

\subsubsection{이번 학습의 핵심 질문}

이번 학습에서 설명할 수 있어야 하는 핵심 질문은 다음과 같다.

\begin{dhksmry}
- def는 함수 본문을 실행하는가, 함수 객체를 생성하는가?

- 함수의 이름과 함수 객체는 어떤 관계인가?

- func와 func()는 각각 무엇을 평가하는가?

- 함수 객체가 아닌데도 ()로 호출할 수 있는 객체는 무엇인가?

- 호출할 때 인수는 어떻게 전달되며 return은 무엇을 반환하는가?

- 이름을 재바인딩하면 기존 함수 객체에도 영향을 주는가?
\end{dhksmry}

이 질문들은 독립적으로 보이지만 모두 하나의 실행 모델로 연결된다.

\begin{dhkgraph}[uniform]
node n1 (1,1) :: def -> function 객체 생성
node n2 (2,1) :: 함수 이름 -> 함수 객체 참조
node n3 (3,1) :: 이름 탐색 -> callable 호출 -> 결과 반환
group g1(n1) :: 1. 함수 정의문 실행
group g2(n2) [accent] :: 2. 이름 바인딩
group g3(n3) :: 3. 호출과 반환

edge g1 -> g2
edge g2 -> g3

caption :: 함수 정의, 이름 바인딩, 호출은 서로 연결되지만 별개로 분석해야 하는 실행 단계다.
\end{dhkgraph}

이제 각 단계를 실제 Python 객체 모델에 대응시켜 보자.

\subsubsection{함수 객체(Function Object)의 정확한 정의}

Python에서 함수(Function)는 단순한 코드 블록이 아니다.

함수 객체(Function Object)는 호출 가능한 객체로, 실행할 코드와 함수 실행에 필요한 메타데이터 및 관련 네임스페이스 정보를 가진 객체다.

다음 코드를 보자.

\begin{dhkcode}
def add(a,b):
    return a + b
\end{dhkcode}

이 코드를 실행하면 Python은 \zsh{add()}의 본문을 즉시 실행하지 않는다.

대신 \zsh{def} 정의문을 실행하면서 새로운 함수 객체를 만들고, 현재 네임스페이스의 \zsh{add}라는 이름을 그 객체에 바인딩한다.

이 과정을 구분하면 다음과 같다.

\begin{dhkgraph}[uniform]
group g1(n1) :: 모듈 네임스페이스
node n1 (1,1) :: add

group g2(n2) [accent] :: 함수 객체
node n2 (1,2) :: type: function // __name__: "add" // __code__: 실행 코드

edge g1 -> g2

explain :: add -> 함수 객체에 대한 참조 // 함수 본문은 아직 실행되지 않음 // 정의문 실행 시점의 구조
\end{dhkgraph}

여기서 구분해야 할 것은 세 가지다.

\begin{dhktable}
|구성 요소|의미|
|:---|:---|
|add|함수 객체를 참조하는 이름(Name)|
|함수 객체|실제로 생성된 function 타입의 객체|
|함수 본문|해당 함수가 호출될 때 실행되는 코드|
\end{dhktable}

함수 객체와 그 객체를 가리키는 이름은 동일한 것이 아니다.

이것은 클래스 객체 모델에서 학습하는 다음 구조와 본질적으로 같다.

\begin{dhkcode}
p = Person("Kim")
\end{dhkcode}

\zsh{p}라는 이름이 \zsh{Person} 인스턴스를 참조하는 것처럼, \zsh{add}라는 이름도 함수 객체를 참조하는 것이다.

다만 생성되는 객체의 타입과 기능이 다르다.

\begin{dhknote}
\begin{dhksmry}[centered]
함수 객체의 타입 확인
\end{dhksmry}
\end{dhknote}

\begin{dhkcode}
from types import FunctionType

def add(a, b):
    return a + b

print(type(add))
print(isinstance(add, FunctionType))
print(callable(add))
\end{dhkcode}

예상 실행 결과:

\begin{dhkzsh}
<class 'function'>
True
True
\end{dhkzsh}

여기에서 확인할 수 있는 사실은 다음과 같다.

\zsh{add}는 이름이면서 표현식으로 사용될 때 함수 객체를 참조한다.

\zsh{type(add)}는 그 객체의 실제 타입을 반환하고, \zsh{callable(add)}는 그 객체가 호출 가능한지를 검사한다.

또한 \zsh{FunctionType}은 Python에서 사용자 정의 함수 객체를 나타내므로,

\zsh{isinstance(add, FunctionType)}이 \zsh{True}가 되는 것이다.

중요한 것은 함수 정의문을 실행한 것만으로도 호출 가능한 객체가 만들어진다는 점이다.

\subsubsection{\zsh{def}는 실행문이다 - 함수 생성과 함수 실행의 구분}

다음 코드를 살펴보자.

\begin{dhkcode}
def add(a, b):
    print("함수 본문 실행")
    return a + b


print("함수 정의 완료")

result = add(3, 5)

print(result)
\end{dhkcode}

예상 실행 결과:

\begin{dhkzsh}
함수 정의 완료
함수 본문 실행
8
\end{dhkzsh}

왜 이런 순서가 발생할까?

\begin{dhknote}
\begin{dhksmry}[centered]
1단계. 함수 정의문 실행
\end{dhksmry}
\end{dhknote}

Python이 다음 문장을 만났다.

\begin{dhkcode}
def add(a, b):
    print("함수 본문 실행")
    return a + b
\end{dhkcode}

이때 함수 객체가 생성되고 \zsh{add} 이름에 바인딩된다.

그러나 본문에 있는 \zsh{print()}와 \zsh{return}은 아직 실행되지 않는다.

\begin{dhknote}
\begin{dhksmry}[centered]
2단계. 다음 문장 실행
\end{dhksmry}
\end{dhknote}

\begin{dhkcode}
print("함수 정의 완료")
\end{dhkcode}

이 문장은 함수 바깥의 모듈 수준 코드다.

따라서 즉시 실행되어 첫 번째 문자열이 출력된다.

\begin{dhknote}
\begin{dhksmry}[centered]
3단계. 함수 객체 호출
\end{dhksmry}
\end{dhknote}

\begin{dhkcode}
result = add(3,5)
\end{dhkcode}

여기서 처음으로 \zsh{add}가 참조하는 함수 객체를 호출한다.

호출 과정에서 3과 5에 해당하는 인수 객체가 각각 매개변수 \zsh{a}와 \zsh{b}에 바인딩되고, 함수 본문이 실행된다.

\begin{dhknote}
\begin{dhksmry}[centered]
4단계. 반환값에 이름 바인딩
\end{dhksmry}
\end{dhknote}

\begin{dhkcode}
return a + b
\end{dhkcode}

정수 3과 5의 덧셈 결과인 8이 반환된다.

그 결과가 \zsh{result}라는 이름에 바인딩된다.

\begin{dhknote}
\begin{dhksmry}[centered]
5단계. 반환된 객체 사용
\end{dhksmry}
\end{dhknote}

\begin{dhkcode}
print(result)
\end{dhkcode}

이제 \zsh{result}는 함수 객체가 아니라 정수 객체 8을 참조하고 있다.

따라서 8이 출력된다.

여기서 얻을 결론은 명확하다.

\zsh{def} 문장이 실행되는 시점과 함수 본문이 실행되는 시점은 다르다.

\begin{dhknote}
\begin{dhksmry}[centered]
그렇다면 def를 실행할 때 아무런 계산도 발생하지 않을까?
\end{dhksmry}
\end{dhknote}

그렇지는 않다.

\zsh{def} 문장은 함수 본문을 실행하지 않을 뿐, 함수 정의에 포함된 다른 표현식은 평가할 수 있다.

대표적인 사례가 기본 인수(Default Argument)다.

\begin{dhkcode}
def create_default():
    print("기본값 생성")
    return []

def record(value, history=create_default()):
    history.append(value)
    return history

print("함수 정의 완료")

print(record(10))
print(record(20))
\end{dhkcode}

예상 실행 결과:

\begin{dhkzsh}
기본값 생성
함수 정의 완료
[10]
[10, 20]
\end{dhkzsh}

여기서는 \zsh{record()}가 호출되기도 전에 \zsh{"기본값 생성"}이 출력된다.

왜냐하면 \zsh{history=create_default()}의 기본값 표현식은 \zsh{record} 함수의 정의문이 실행될 때 평가되기 때문이다.

그리고 평가 결과인 리스트 객체 하나가 함수의 기본 인수로 유지된다.

따라서 기본 인수를 생략한 두 호출은 같은 리스트 객체를 사용하고, 첫 번째 호출의 변경이 두 번째 호출에서도 관찰되는 것이다.

이 원리는 Python 공식 언어 참조에도 명시되어 있다.

따라서 다음 두 명제는 구별해야 한다.

\begin{dhktable}
|명제|판단|
|:---|:---|
|함수 정의문 실행은 함수 본문을 실행하지 않는다.|참|
|함수 정의문 실행 중에는 어떤 표현식도 평가하지 않는다.|거짓|
\end{dhktable}

데코레이터 표현식 역시 함수 정의 시점에 평가될 수 있다.

다만 그 평가·적용 순서는 이후에 별도로 분석할 것이다.

\subsubsection{이름 바인딩 - 하나의 함수 객체를 여러 이름이 참조할 수 있는가?}

가능하다.

그리고 이것이 일급 함수(First-Class Function)를 이해하기 위한 첫 번째 조건이다.

다음 코드를 보자.

\begin{dhkcode}
def add(a, b):
    return a + b

operation = add

print(operation is add)
print(operation(3, 5))
\end{dhkcode}

예상 실행 결과:

\begin{dhkzsh}
True
8
\end{dhkzsh}

여기서 중요한 문장은 다음 문장이다.

\begin{dhkcode}
operation = add
\end{dhkcode}

이 문장은 \zsh{add} 함수 객체를 복사하거나 새로 생성하지 않는다.

\zsh{add}가 참조하는 기존 함수 객체를 \zsh{operation}이라는 이름도 참조하도록 바인딩하는 것이다.

\begin{dhkgraph}[uniform]
group g1(n1, n2) :: 네임스페이스
node n1 (1,1) :: add
node n2 (2,1) :: operation
group g2(n3) :: 함수 객체
node n3 (2,2) :: function

edge n1 -> g2
edge n2 -> g2

explain :: operation is add -> True // 두 이름이 동일한 객체를 참조함
\end{dhkgraph}

\begin{dhknote}
\begin{dhksmry}[centered]
이름이 바뀌면 함수 객체도 바뀔까?
\end{dhksmry}
\end{dhknote}

이번에는 다음 코드를 살펴보자.

\begin{dhkcode}
def double(x):
    return x * 2

saved = double

double = "replaced"

print(callable(double))
print(callable(saved))
print(saved(5))
print(saved.__name__)
\end{dhkcode}

예상 실행 결과:

\begin{dhkzsh}
False
True
10
double
\end{dhkzsh}

이 과정은 객체 생성, 참조, 이름 재바인딩으로 나눠서 분석할 수 있다.

\begin{dhknote}
\begin{dhksmry}[centered]
첫 번째 단계 - 함수 객체 생성
\end{dhksmry}
\end{dhknote}

\begin{dhkcode}
def double(x):
    return x * 2
\end{dhkcode}

함수 객체가 생성되고 \zsh{double}이라는 이름에 바인딩된다.

\begin{dhknote}
\begin{dhksmry}[centered]
두 번째 단계 - 참조 공유
\end{dhksmry}
\end{dhknote}

\begin{dhkcode}
saved = double
\end{dhkcode}

\zsh{saved}도 동일한 함수 객체를 참조하게 된다.

\begin{dhknote}
\begin{dhksmry}[centered]
세 번째 단계 - 이름 재바인딩
\end{dhksmry}
\end{dhknote}

\begin{dhkcode}
double = "replaced"
\end{dhkcode}

이 문장은 원래의 함수 객체를 수정하는 것이 아니다.

\zsh{double}이라는 이름을 새로운 문자열 객체에 다시 바인딩한다.

기존 함수 객체는 여전히 \zsh{saved}가 참조하고 있다.

\begin{dhknote}
\begin{dhksmry}[centered]
네 번째 단계 - 기존 함수 객체 호출
\end{dhksmry}
\end{dhknote}

\begin{dhkcode}
saved(5)
\end{dhkcode}

\zsh{saved}는 원래의 함수 객체를 가리키므로 정상적으로 실행된다.

또한 \zsh{saved.__name__}이 \zsh{"double"}이라는 사실도 중요하다.

\zsh{__name__}은 함수 객체에 있는 메타데이터다.

함수 객체를 참조하는 현재 변수 이름과 반드시 일치할 필요는 없다.

따라서 다음 두 가지를 구분해야 한다.

\begin{dhksmry}
    - 이름 재바인딩(Rebinding) : double = "replaced"

    - 객체 상태 변경(Mutation) : 객체 자체가 가진 정보를 변경
\end{dhksmry}

위 예제는 이름 재바인딩이지 함수 객체의 변경이 아니다.

이 구조는 나중에 다음 데코레이터 변환식을 이해하는 핵심이 된다.

\begin{dhkcode}
add = decorator(add)
\end{dhkcode}

여기서 \zsh{add}라는 이름에 새로운 반환 객체를 바인딩하더라도, 기존 함수 객체를 참조하는 다른 이름이나 클로저가 있다면 기존 함수 객체는 계속 사용될 수 있다.

\subsubsection{\zsh{func}와 \zsh{func()} - 함수 참조와 호출의 차이}

이번에는 호출 표현식(Call Expression)을 살펴보자.

\begin{dhkcode}
def greet():
    print("Hello")
    return 100
\end{dhkcode}

다음 두 표현식은 의미가 다르다.

\begin{dhkcode}
greet
\end{dhkcode}

이것은 이름을 탐색하여 함수 객체를 얻는 표현식이다.

반면

\begin{dhkcode}
greet()
\end{dhkcode}

이것은 \zsh{greet}이 참조하는 호출 가능한 객체를 호출하는 표현식이다.

실제 코드를 비교해 보자.

\begin{dhkcode}
def greet():
    print("Hello")
    return 100

a = greet
b = greet()

print(type(a).__name__)
print(b)
\end{dhkcode}

예상 실행 결과:

\begin{dhkzsh}
Hello
function
100
\end{dhkzsh}

\zsh{a = greet}에서는 함수 본문이 실행되지 않는다.

그 대신 \zsh{a}가 함수 객체를 참조하게 된다.

반면 \zsh{b = greet()}에서는 함수 객체를 호출하여 \zsh{"Hello"}를 출력하고, 100이라는 정수 객체를 반환받아 \zsh{b}에 바인딩한다.

\begin{dhknote}
\begin{dhksmry}[centered]
()의 의미
\end{dhksmry}
\end{dhknote}

여기에서 하나 더 정확하게 이해할 부분이 있다.

\zsh{()}는 단순히 함수만 실행하는 문법이 아니다.

\zsh{()}는 앞의 표현식을 평가하여 얻은 호출 가능한 객체(Callable)를 호출하는 문법이다.

따라서 다음 표현식들은 모두 호출 표현식이다.

\begin{dhkcode}
func()
obj.method()
SomeClass()
callable_instance()
\end{dhkcode}

각각 호출하는 객체의 종류는 다를 수 있다.

\begin{dhktable}
|표현식|평가하여 호출하는 대상|
|:---|:---|
|func()|일반적으로 함수 객체|
|obj.method()|일반적으로 바운드 메서드 객체|
|SomeClass()|클래스 객체|
|callable_instance()|__call__을 지원하는 인스턴스|
\end{dhktable}

또한 함수 이름으로 시작하지 않는 표현식도 호출할 수 있다.

\begin{dhkcode}
def get_function():
    return print

get_function()("Hello")
\end{dhkcode}

예상 실행 결과:

\begin{dhkzsh}
Hello
\end{dhkzsh}

실행 흐름을 따라가 보자.

먼저 \zsh{get_function()}이 호출된다.

이 함수는 \zsh{print}가 참조하는 내장 함수 객체를 반환한다.

그러면 바깥쪽 호출 표현식은 그 반환된 함수 객체를 \zsh{"Hello"}라는 인수로 다시 호출한다.

개념적으로 다음 과정과 같다.

\begin{dhkcode}
f = get_function()
f("Hello")
\end{dhkcode}

이 원리는 추후 인수를 받는 데코레이터의 다음 구조를 분석할 때 다시 등장한다.

\begin{dhkcode}
decorator_factory(config)(func)
\end{dhkcode}

여기서 첫 번째 호출의 반환값이 또 다른 호출 가능한 객체라는 사실이 중요하다.

\subsubsection{\zsh{callable()} - 호출 가능한 객체는 함수뿐일까?}

아니다.

Python은 함수 객체뿐 아니라 클래스 객체, 메서드 객체, 호출 가능하도록 설계된 인스턴스도 호출할 수 있다.

\zsh{callable(obj)}는 어떤 객체가 호출 가능한지 검사하는 내장 함수다.

다만 \zsh{True}가 나왔다고 해서 모든 인수 조합으로 호출이 성공한다는 뜻은 아니다.

\begin{dhknote}
\begin{dhksmry}[centered]
사용자 정의 인스턴스를 호출 가능하게 만들기
\end{dhksmry}
\end{dhknote}

클래스 객체 모델에서 살펴보는 \zsh{__call__()}을 연결해 보자.

\begin{dhkcode}
class Multiplier:
    def __init__(self, factor):
        self.factor = factor

    def __call__(self, value):
        return self.factor * value

multiply = Multiplier(3)

print(type(multiply).__name__)
print(callable(multiply))
print(multiply(4))
print(multiply.__call__.__self__ is multiply)
\end{dhkcode}

예상 실행 결과:

\begin{dhkzsh}
Multiplier
True
12
True
\end{dhkzsh}

여기서 \zsh{multiply}는 함수 객체가 아니다.

\zsh{Multiplier} 인스턴스다.

그런데도 \zsh{multiply(4)}라는 표현식이 가능하다.

왜 그럴까?

\begin{dhknote}
\begin{dhksmry}[centered]
호출 과정 분석
\end{dhksmry}
\end{dhknote}

먼저 다음 표현식을 실행한다.

\begin{dhkcode}
multiply = Multiplier(3)
\end{dhkcode}

클래스 객체 \zsh{Multiplier}가 호출되고 새로운 인스턴스가 생성된다.

\zsh{__init__()}이 실행되면서 해당 인스턴스의 \zsh{factor} 속성에 3이 바인딩된다.

다음 표현식은

\begin{dhkcode}
multiply(4)
\end{dhkcode}

\zsh{multiply} 인스턴스를 호출하는 표현식이다.

Python의 특수 메서드 호출 메커니즘은 그 인스턴스의 타입에서 \zsh{__call__} 구현을 찾아 사용한다.

이 예제에서는 개념적으로 다음과 같은 효과를 가진다.

\begin{dhkcode}
Multiplier.__call__(multiply, 4)
\end{dhkcode}

\zsh{self}는 \zsh{multiply} 인스턴스를 참조하고, \zsh{value}는 정수 4를 참조한다.

따라서 다음 연산이 수행된다.

\begin{dhkcode}
self.factor * value
\end{dhkcode}

결과는 \zsh{3 * 4}, 즉 12이다.

일반적인 사용자 정의 클래스의 특수 메서드는 인스턴스 딕셔너리에서 단순히 같은 이름을 찾는 방식과 달리, 타입을 통한 특별한 탐색 규칙이 적용된다는 점도 기억한다.

\begin{dhknote}
\begin{dhksmry}[centered]
중요한 비교
\end{dhksmry}
\end{dhknote}

\begin{dhktable}
|구분|	add|	multiply|
|:---|:---|:---|
|객체 타입|	function|	Multiplier|
|callable()|	True|	True|
|호출 표현식|	add(3, 4)|	multiply(4)|
|호출 동작|	함수 본문 실행|	클래스에 정의된 __call__ 실행|
|상태 보관|	함수 속성·기본값·클로저 등|	일반적으로 인스턴스 속성|
\end{dhktable}

이로부터 다음 관계가 성립한다.

모든 일반적인 사용자 정의 함수 객체는 호출 가능하지만, 모든 호출 가능한 객체가 함수 객체인 것은 아니다.

여기서 기존 클래스 객체 모델과 연결되는 중요한 통찰이 하나 있다.

PyTorch의 \zsh{model(x)} 역시 함수 객체를 호출하는 것이 아니다.

\zsh{nn.Module}을 상속한 모델 인스턴스를 호출하는 표현식이다.

반면 \zsh{model.forward(x)}는 속성 탐색을 통해 얻은 바운드 메서드를 호출한다.

두 표현식 모두 \zsh{()}를 사용하지만 호출 대상과 실행 경로는 다르다.

\begin{dhknote}
\begin{dhksmry}[centered]
callable()이 True여도 오류가 발생하는 경우
\end{dhksmry}
\end{dhknote}

\begin{dhkcode}
def add(a, b):
    return a + b

print(callable(add))

add(10)
\end{dhkcode}

첫 번째 출력은 다음과 같다.

\begin{dhkzsh}
True
\end{dhkzsh}

하지만 두 번째 호출에서는 \zsh{TypeError}가 발생한다.

\zsh{add}는 분명히 호출 가능한 함수 객체이지만 필수 매개변수 \zsh{b}에 대응하는 인수를 전달하지 않았기 때문이다.

이 사례는 세 가지 판단을 분리해야 한다는 의미다.

\begin{dhktable}
|판단|질문|
|:---|:---|
|호출 가능성|이 객체를 호출할 수 있는가?|
|인수 적합성|전달한 인수가 함수의 매개변수 규약에 맞는가?|
|실행 성공|함수 본문이 예외 없이 정상 종료되는가?|
\end{dhktable}

\zsh{callable()}은 첫 번째 판단을 위한 기능이다.

나머지 두 가지까지 보장하지는 않는다.

\subsubsection{함수 호출 시 인수는 어떻게 전달되는가?}

클래스 객체 모델에서는 인스턴스를 생성하거나 메서드를 호출할 때 인수를 전달한다.

\begin{dhkcode}
Person("Kim")
p.introduce()
\end{dhkcode}

이제 이 원리를 함수 호출에 적용해 보자.

다음 코드를 살펴보자.

\begin{dhkcode}
def update(values):
    values.append("B")

    values = ["C"]

    return values

original = ["A"]

result = update(original)

print(original)
print(result)
print(original is result)
\end{dhkcode}

예상 실행 결과:

\begin{dhkzsh}
['A', 'B']
['C']
False
\end{dhkzsh}

왜 이 결과가 발생할까?

\begin{dhknote}
\begin{dhksmry}[centered]
호출 전 상태
\end{dhksmry}
\end{dhknote}

\begin{dhkcode}
original = ["A"]
\end{dhkcode}

리스트 객체 하나가 생성되고 \zsh{original}이라는 이름에 바인딩된다.

이후 다음 표현식이 실행된다.

\begin{dhkcode}
update(original)
\end{dhkcode}

Python은 \zsh{original}을 평가하여 리스트 객체를 얻고, 함수의 매개변수 \zsh{values}에 같은 객체를 바인딩한다.

이 시점의 구조는 다음과 같다.

\begin{dhkgraph}[uniform]
group g1(n1) :: 이름이 존재하는 공간
node n1 (1,1) :: original(모듈) // values(함수 지역)
group g2(n2) [accent] :: 리스트 객체
node n2 (1,2) :: ["A"]

edge g1 -- g2

explain :: 두 이름이 동일한 리스트 객체를 참조
\end{dhkgraph}

이후 함수 본문에서 다음 문장이 실행된다.

\begin{dhkcode}
values.append("B")
\end{dhkcode}

\zsh{values}가 참조하는 리스트 객체 자체의 상태가 변경된다.

따라서 \zsh{original}이 참조하는 리스트에도 \zsh{"B"}가 추가된 것이 관찰되는 것이다.

다음 문장은 다르게 동작한다.

\begin{dhkcode}
values = ["C"]
\end{dhkcode}

이 문장은 새로운 리스트 객체를 만들고 함수의 지역 이름 \zsh{values}를 그 객체에 다시 바인딩한다.

기존 리스트 객체는 여전히 \zsh{original}이 참조하고 있다.

마지막으로

\begin{dhkcode}
return values
\end{dhkcode}

새로운 리스트 객체가 반환되고, 호출한 쪽의 \zsh{result} 이름에 바인딩된다.

따라서 최종 관계는 다음과 같다.

\begin{dhksmry}
    original \ra ["A", "B"]

    result \ra ["C"]
\end{dhksmry}

이 예제에서 반드시 구분해야 할 연산은 다음 세 가지다.

\begin{dhktable}
|표현식|발생한 연산|
|:---|:---|
|update(original)|인수 객체를 매개변수 이름에 바인딩|
|values.append("B")|기존 리스트 객체의 상태 변경|
|values = ["C"]|새로운 리스트 객체 생성 및 지역 이름 바인딩|
\end{dhktable}

이 원리는 데코레이터의 Wrapper에서도 동일하게 적용된다.

예를 들어 나중에 다음 코드를 보면

\begin{dhkcode}
def wrapper(*args, **kwargs):
    return func(*args, **kwargs)
\end{dhkcode}

\zsh{args}와 \zsh{kwargs}가 어떤 객체를 담고 있고, 그 객체들이 원본 함수의 매개변수에 어떻게 전달되는지 분석할 수 있어야 한다.

\subsubsection{\zsh{return} - 무엇을 반환하는가?}

함수는 반드시 숫자나 문자열을 반환해야 하는 것이 아니다.

\zsh{return}은 함수 실행을 종료하고, 지정된 표현식을 평가하여 얻은 객체를 호출한 쪽에 반환하는 문장이다.

\begin{dhknote}
\begin{dhksmry}[centered]
첫 번째 경우 - 정수 객체 반환
\end{dhksmry}
\end{dhknote}

\begin{dhkcode}
def compute():
    return 10

result = compute()
\end{dhkcode}

\zsh{compute()}가 호출되고 정수 객체 10이 반환된다.

그 결과 \zsh{result}가 정수 객체를 참조하게 된다.

\begin{dhknote}
\begin{dhksmry}[centered]
두 번째 경우 - 함수 객체 반환
\end{dhksmry}
\end{dhknote}

\begin{dhkcode}
def double(x):
    return x * 2

def get_operation():
    return double

operation = get_operation()

print(operation is double)
print(operation(5))
\end{dhkcode}

예상 실행 결과:

\begin{dhkzsh}
True
10
\end{dhkzsh}

이 코드에서 \zsh{get_operation()}은 정수나 문자열이 아니라 함수 객체를 반환한다.

따라서 반환된 함수 객체를 \zsh{operation}이라는 이름에 바인딩한 다음, 호출할 수 있다.

여기서 \zsh{get_operation()} 자체는 중첩 함수를 정의하지 않았다.

이미 존재하는 함수 객체를 반환했을 뿐이다.

그러므로 이 예제에는 클로저가 필요하지 않다.

\begin{dhknote}
\begin{dhksmry}[centered]
세 번째 경우 - 함수 호출 결과 반환
\end{dhksmry}
\end{dhknote}

이번에는 다음 코드를 살펴보자.

\begin{dhkcode}
def double(x):
    return x * 2

def run_operation():
    return double(5)

result = run_operation()

print(result)
print(type(result).__name__)
\end{dhkcode}

예상 실행 결과:

\begin{dhkzsh}
10
int
\end{dhkzsh}

이번에는 \zsh{double}이라는 함수 객체를 반환하는 것이 아니라, \zsh{double(5)}를 먼저 호출하고, 그 실행 결과인 정수 객체를 반환한다.

\begin{dhknote}
\begin{dhksmry}[centered]
return func와 return func() 비교
\end{dhksmry}
\end{dhknote}

\begin{dhktable}
|표현식|실행 의미|반환 대상|
|:---|:---|:---|
|return func|함수 객체를 평가한 뒤 반환|함수 객체|
|return func()|함수 객체를 호출한 뒤 결과 반환|호출 결과 객체|
\end{dhktable}

이 차이는 \zsh{A.6}과 \zsh{A.7}에서 매우 중요해진다.

예를 들어 Wrapper 함수 객체를 반환해야 하는 데코레이터에서

\begin{dhkcode}
return wrapper
\end{dhkcode}

대신

\begin{dhkcode}
return wrapper()
\end{dhkcode}

를 사용하면 어떻게 될까?

후자의 경우 Wrapper를 반환하기 전에 즉시 호출하려고 한다.

인수가 필요한 Wrapper라면 \zsh{TypeError}가 발생할 수도 있다.

인수가 필요 없더라도 Wrapper의 실행 결과가 반환되므로 의도했던 함수 객체가 반환되지 않을 수 있다.

물론 Wrapper의 호출 결과 자체가 다시 호출 가능한 객체라면 이후 호출이 가능할 수도 있지만, 그것은 별도로 설계한 동작이지 일반적인 데코레이터의 기본 구조는 아니다.

\begin{dhknote}
\begin{dhksmry}[centered]
return이 없으면 무엇이 반환되는가?
\end{dhksmry}
\end{dhknote}

다음 코드를 보자.

\begin{dhkcode}
def show(value):
    print(value)

result = show(10)

print(result is None)
\end{dhkcode}

예상 실행 결과:

\begin{dhkzsh}
10
True
\end{dhkzsh}

\zsh{show()}에는 명시적인 \zsh{return}이 없다.

함수 본문이 정상적으로 끝나면 Python은 \zsh{None}을 반환한다.

여기서 \zsh{print(value)}가 화면에 10을 출력했다는 사실과 \zsh{show()}가 10을 반환했다는 것은 서로 다른 이야기이다.

\zsh{print()}의 출력은 관찰 가능한 부작용(Side Effect)이고, 함수의 반환값은 호출 표현식이 평가된 결과다.

또한 함수 안에서 인수 없는 \zsh{return}을 실행하는 경우에도 \zsh{None}이 반환된다.

이는 일반적인 \zsh{def} 함수의 실행을 설명한 것이며, 제너레이터 함수나 코루틴 함수에는 추가적인 호출·실행 규칙이 적용된다.

그 차이는 \zsh{A.13}에서 따로 다룰 것이다.

\subsubsection{일급 함수(First-Class Function) - 함수를 데이터처럼 다루기}

지금까지 함수 객체를 다른 이름에 바인딩할 수도 있고, 반환할 수도 있다는 사실을 확인했다.

이러한 성질은 Python에서 함수가 일급 객체(First-Class Object)로 취급되기 때문에 가능하다.

일급 객체라는 표현은 특정 객체를 일반적인 값처럼 변수에 바인딩하고, 자료구조에 저장하고, 함수에 전달하고, 반환할 수 있다는 의미로 사용된다.

다음 예제를 보자.

\begin{dhkcode}
def add(a, b):
    return a + b

def multiply(a, b):
    return a * b

def apply(operation, a, b):
    return operation(a, b)

operations = {
    "sum": add,
    "product": multiply,
}

print(apply(operations["sum"], 3, 4))
print(apply(operations["product"], 3, 4))
\end{dhkcode}

예상 실행 결과:

\begin{dhkzsh}
7
12
\end{dhkzsh}

이 코드는 단순하지만 중요한 구조를 포함하고 있다.

\begin{dhknote}
\begin{dhksmry}[centered]
객체 생성 및 저장
\end{dhksmry}
\end{dhknote}

먼저 \zsh{add}, \zsh{multiply}, \zsh{apply}라는 세 개의 함수 객체가 생성된다.

그리고 다음 딕셔너리를 만든다.

\begin{dhkcode}
operations = {
    "sum": add,
    "product": multiply,
}
\end{dhkcode}

이 딕셔너리의 값은 함수 실행 결과가 아니다.

함수 객체에 대한 참조다.

따라서 다음 구조가 형성된다.

\begin{dhkdir}[in=6]
operations
> "sum" -> add 함수 객체
> "product" -> multiply 함수 객체
\end{dhkdir}

\begin{dhknote}
\begin{dhksmry}[centered]
함수 객체를 인수로 전달
\end{dhksmry}
\end{dhknote}

다음 표현식을 실행하면

\begin{dhkcode}
apply(operations["sum"], 3, 4)
\end{dhkcode}

먼저 딕셔너리에서 \zsh{"sum"}에 대응하는 함수 객체를 조회한다.

그 함수 객체가 \zsh{apply}의 첫 번째 매개변수 \zsh{operation}에 전달된다.

그리고 3, 4가 각각 \zsh{a}, \zsh{b}에 전달된다.

함수 본문에서는 다음 호출이 실행된다.

\begin{dhkcode}
operation(a, b)
\end{dhkcode}

\zsh{operation}은 이 호출에서 \zsh{add} 함수 객체를 참조하고 있으므로 \zsh{add(3, 4)}가 실행되는 것이다.

\begin{dhknote}
\begin{dhksmry}[centered]
고차 함수(Higher-Order Function)와의 관계
\end{dhksmry}
\end{dhknote}

여기서 \zsh{apply}는 고차 함수다.

고차 함수는 다른 함수를 인수로 받거나 함수를 반환하는 함수를 의미한다.

따라서 일급 함수와 고차 함수는 관계가 있지만 동일한 개념은 아니다.

일급 함수는 함수 객체를 어떤 값으로 취급할 수 있는가에 관한 성질이고, 고차 함수는 그 함수 객체를 인수나 반환값으로 사용하는 함수의 형태를 설명하는 개념이다.

이 구분은 데코레이터를 학습할 때 다시 연결된다.

데코레이터 함수는 보통 다른 함수 객체를 인수로 받으므로 고차 함수의 한 사례가 될 수 있다.

\subsubsection{함수 객체의 내부 정보 - \zsh{__name__}, \zsh{__code__}, \zsh{__globals__}}

이제 함수가 정말 하나의 객체라면 어떤 속성을 가지고 있는지도 확인해 보자.

다음 코드를 살펴보자.

\begin{dhkcode}
def add(a, b=10):
    return a + b

print(add.__name__)
print(add.__defaults__)
print(add.__globals__ is globals())
print(add.__closure__)
print(type(add.__code__).__name__)
\end{dhkcode}

예상 실행 결과:

\begin{dhkzsh}
add
(10,)
True
None
code
\end{dhkzsh}

이 결과는 함수 객체가 단순히 실행 코드 하나만 가진 구조가 아니라는 사실을 보여준다.

\begin{dhktable}
|속성|의미|
|:---|:---|
|__name__|	함수의 이름을 나타내는 메타데이터|
|__qualname__|	함수가 정의된 렉시컬 구조를 반영하는 정규화된 이름|
|__code__|	함수의 실행 코드 정보를 담는 코드 객체|
|__globals__|	함수가 정의된 모듈의 전역 네임스페이스|
|__defaults__|	위치 매개변수 기본값을 담는 튜플 또는 None|
|__kwdefaults__|	키워드 전용 매개변수 기본값 정보|
|__closure__|	캡처한 자유 변수와 연결된 셀 객체들의 튜플 또는 None|
|__dict__|	일반 사용자 정의 함수 객체에 추가한 속성을 저장하는 딕셔너리|
\end{dhktable}

이러한 속성은 Python의 함수 데이터 모델에 정의되어 있다.

\begin{dhknote}
\begin{dhksmry}[centered]
__code__와 함수 객체는 동일한가?
\end{dhksmry}
\end{dhknote}

아니다.

\zsh{__code__}는 코드 객체(Code Object)를 가리킨다.

반면 함수 객체는 코드 객체 이외에도 기본 인수, 네임스페이스 참조, 클로저 등의 정보를 가지고 있다.

따라서 함수 객체를 코드 객체 하나로 환원해서 이해하면 안 된다.

\begin{dhknote}
\begin{dhksmry}[centered]
__globals__는 무엇을 참조하는가?
\end{dhksmry}
\end{dhknote}

다음 코드를 보자.

\begin{dhkcode}
RATE = 2

def calculate(x):
    return x * RATE

print(calculate(4))

RATE = 3

print(calculate(4))
\end{dhkcode}

예상 실행 결과:

\begin{dhkzsh}
8
12
\end{dhkzsh}

여기서 \zsh{calculate} 함수 객체는 처음 정의되었을 때 전역 변수 \zsh{RATE}의 값 2를 복사해 저장한 것이 아니다.

함수가 정의된 모듈의 전역 네임스페이스를 참조하고 있고, 호출 시 \zsh{RATE}를 조회할 때 그 네임스페이스에 현재 바인딩된 객체를 사용하는 것이다.

따라서 \zsh{RATE = 3}으로 전역 이름을 재바인딩한 뒤 함수를 다시 호출하면 결과가 달라진다.

이것은 \zsh{A.3}의 이름 탐색과 연결되는 중요한 예제다.

\begin{dhknote}
\begin{dhksmry}[centered]
__globals__와 __closure__의 차이
\end{dhksmry}
\end{dhknote}

위 예제의 \zsh{RATE}는 전역 이름이다.

따라서 이 예제에서는 \zsh{calculate.__closure__}가 \zsh{None}이더라도 \zsh{RATE}에 접근할 수 있다.

클로저가 외부 함수의 지역 바인딩을 어떻게 보존하는지는 \zsh{A.4}에서 학습할 것이다.

이때는 \zsh{__closure__}, 코드 객체의 \zsh{co_freevars}, 셀 객체(Cell Object)를 연결해 살펴본다.

\subsubsection{클래스 객체 모델과 다시 연결하기}

이제 \zsh{A.1}에서 학습한 함수 객체를 클래스 내부에 배치해 보자.

\begin{dhkcode}
class Calculator:
    def add(self, a, b):
        return a + b

calc = Calculator()

func = Calculator.__dict__["add"]

method = calc.add

print(type(func).__name__)
print(type(method).__name__)

print(method.__func__ is func)
print(method.__self__ is calc)

print(func(calc, 3, 4))
print(method(3, 4))
\end{dhkcode}

예상 실행 결과:

\begin{dhkzsh}
function
method
True
True
7
7
\end{dhkzsh}

여기에서는 같은 함수를 서로 다른 경로로 호출하고 있다.

첫 번째 경로를 살펴보자.

\begin{dhkcode}
func = Calculator.__dict__["add"]
\end{dhkcode}

클래스 네임스페이스에서 \zsh{add}에 대응하는 함수 객체를 직접 얻는다.

따라서 \zsh{func}의 타입은 \zsh{function}이다.

다음 호출은

\begin{dhkcode}
func(calc, 3, 4)
\end{dhkcode}

함수 객체를 직접 호출하면서 첫 번째 인수로 \zsh{calc}를 전달하는 구조다.

반면 두 번째 경로에서는

\begin{dhkcode}
method = calc.add
\end{dhkcode}

\zsh{calc.add} 속성에 접근할 때 함수 객체의 디스크립터 메커니즘을 통해 \zsh{calc}에 바인딩된 메서드 객체를 얻는다.

따라서 다음 호출에서는

\begin{dhkcode}
method(3, 4)
\end{dhkcode}

이미 바인딩된 \zsh{calc}가 메서드 호출 과정에서 첫 번째 인수로 전달되는 것이다.

이 구조를 비교하면 다음과 같다.

\begin{dhktable}
|구분|func|method|
|:---|:---|:---|
|객체 타입|function|method|
|참조 대상|클래스에 저장된 함수 객체|인스턴스에 바인딩된 메서드 객체|
|self 전달|명시적으로 전달|바인딩된 인스턴스가 전달|
|호출|func(calc, 3, 4)|method(3, 4)|
\end{dhktable}

이렇게 보면 \zsh{A.1}의 함수 객체는 클래스 객체 모델과 분리된 별개의 시스템이 아니다.

클래스 내부의 일반적인 인스턴스 메서드 역시 함수 객체를 바탕으로 디스크립터와 메서드 바인딩 규칙이 결합된 결과라는 사실을 확인할 수 있다.

이 원리는 나중에 데코레이터로 함수 객체를 교체하거나, Wrapper를 클래스 메서드로 사용할 때 특히 중요해진다.

\subsubsection{실제 응용 - 함수 객체를 이용한 명령 처리기}

함수 객체를 이름에 바인딩하고, 딕셔너리에 저장하며, 나중에 호출하는 구조는 실제 프로그램에서도 유용하다.

예를 들어 다음과 같은 간단한 명령 처리 시스템을 만들 수 있다.

\begin{dhkcode}
def greet(name):
    return f"Hello, {name}"

def farewell(name):
    return f"Goodbye, {name}"

handlers = {}

handlers["greet"] = greet
handlers["farewell"] = farewell

def dispatch(command, name):
    handler = handlers[command]

    return handler(name)

print(dispatch("greet", "Kim"))
print(dispatch("farewell", "Lee"))
\end{dhkcode}

예상 실행 결과:

\begin{dhkzsh}
Hello, Kim
Goodbye, Lee
\end{dhkzsh}

이 프로그램에서는 명령 문자열에 따라 다른 함수 객체를 선택해 호출한다.

구조적으로 중요한 부분은 다음이다.

\begin{dhkcode}
handlers["greet"] = greet
\end{dhkcode}

여기서는 \zsh{greet()}을 실행하지 않는다.

대신 \zsh{greet} 함수 객체를 딕셔너리에 등록해 둔다.

나중에 다음 호출이 발생하면

\begin{dhkcode}
dispatch("greet", "Kim")
\end{dhkcode}

딕셔너리에 등록된 함수 객체를 얻는다.

\begin{dhkcode}
handler = handlers[command]
\end{dhkcode}

그 다음에 해당 함수 객체를 호출한다.

\begin{dhkcode}
return handler(name)
\end{dhkcode}

즉, 함수의 등록 시점과 실행 시점이 분리되어 있다.

이것이 추후 등록형 데코레이터를 이해하기 위한 첫 번째 기반이다.

\zsh{A.12}에서는 이러한 명령 처리 구조를 더 발전시켜, \zsh{self.routes} 같은 인스턴스 저장소에 함수 객체를 등록하고 요청 시 호출하는 메커니즘을 직접 구성한 뒤 FastAPI의 경로 등록 구조와 비교할 것이다.

그때는 기존에 배운 클래스 객체, \zsh{self}, 메서드 바인딩과 이번 과정의 클로저까지 모두 연결할 수 있게 된다.

\subsubsection{경계 조건과 흔한 오해}

이번 단원에서 반드시 구별해야 할 표현식들을 표로 정리해 보자.

\begin{dhktable}
|표현식|	정확한 의미|
|:---|:---|
|def f(...): ...|	함수 정의문을 실행하여 함수 객체를 생성하고 이름을 바인딩|
|f|	이름을 탐색하여 현재 바인딩된 객체를 얻음|
|f()|	f가 참조하는 객체를 호출하고 결과를 얻음|
|g = f|	기존 함수 객체에 새로운 이름을 바인딩|
|g = f()|	함수 호출의 반환 객체를 g에 바인딩|
|return f|	함수 객체를 반환|
|return f()|	함수 객체를 호출한 결과를 반환|
|callable(f)|	해당 객체가 호출 가능한지 검사|
|f.__name__|	함수 객체의 이름 메타데이터 조회|
|f.__globals__|	함수가 참조하는 전역 네임스페이스 조회|
|f.__closure__|	캡처된 자유 변수의 셀 객체 정보 조회|
|type(f)|	f가 참조하는 객체의 실제 타입 조회|
\end{dhktable}

한 가지 더 기억할 점이 있다.

함수 객체를 호출하는 행위가 항상 정상적인 반환으로 이어지는 것은 아니다.

호출 과정에서 인수 바인딩 오류가 발생할 수도 있고, 본문 실행 중 예외가 전파될 수도 있으며, 함수가 다른 호출 가능한 객체를 반환할 수도 있다.

또한 함수 객체가 반환되었다고 해서 그것이 반드시 클로저라는 의미는 아니다.

우리가 \zsh{A.1}에서 살펴본 \zsh{get_operation()}의 반환 함수는 외부 함수의 지역 상태를 캡처하지 않았다.

이러한 차이를 유지해야 \zsh{A.2}에서 중첩 함수를 학습하더라도 함수 반환과 클로저 형성을 혼동하지 않게 된다.

\subsubsection{핵심 정리와 판단 기준}

이번 단원의 핵심은 다음 실행 모델이다.

\begin{dhkgraph}[uniform]
node n1 (1,1) :: 함수 정의문 실행 // def f(...): ...

node n2 (2,1) :: 함수 객체 생성 //code · globals · defaults · closure 등

node n3 (3,1):: 이름 바인딩 // f -> function object

node n4 (4,1) :: 호출 가능 여부 및 호출 // callable(f) · f(args)

node n5 (5,1) :: 인수 바인딩 및 함수 본문 실행 // 매개변수 이름 -> 인수 객체

node n6 (6,1) :: 결과 반환 // return expression

edge n1 -> n2
edge n2 -> n3
edge n3 -> n4
edge n4 -> n5
edge n5 -> n6
\end{dhkgraph}

이제 코드를 분석할 때는 다음 판단을 구분하면 된다.

첫째, 새로운 함수 객체가 생성되는 것인지, 이미 존재하는 함수 객체를 참조하는 것인지 판단해야 한다.

둘째, 이름에 객체를 바인딩하는 것인지, 기존 이름을 재바인딩하는 것인지 구분해야 한다.

셋째, \zsh{()}가 호출하는 대상이 함수 객체인지, 클래스 객체인지, 바운드 메서드인지, 아니면 다른 호출 가능한 인스턴스인지 확인해야 한다.

넷째, 호출 시 어떤 객체가 매개변수에 바인딩되고, 함수 본문에서 실제로 어떤 객체가 변경되는지 추적해야 한다.

다섯째, 함수 객체 자체를 반환하는 것과 함수 호출 결과를 반환하는 것을 구별해야 한다.

이 다섯 가지가 \zsh{A.1}의 판단 기준이다.